\chapter*{KHAI BÁO SỬ DỤNG TRỢ LÝ TRÍ TUỆ NHÂN TẠO}
\addcontentsline{toc}{chapter}{Khai báo sử dụng trợ lý Trí tuệ nhân tạo (AI Usage Statement)}

Nhóm sinh viên thực hiện đề tài xin cam kết tuân thủ nghiêm ngặt chuẩn mực về \textbf{Liêm chính học thuật (Academic Integrity)}. Quá trình thực hiện dự án có sử dụng công cụ trợ lý AI (Google Antigravity IDE / Gemini Models) với phạm vi minh bạch như sau:

\begin{enumerate}
    \item \textbf{Phạm vi hỗ trợ của AI:}
    \begin{itemize}
        \item Khởi tạo cấu trúc khung dự án phần mềm theo mô hình Machine Learning chuẩn mực.
        \item Gợi ý cú pháp định dạng mã nguồn Python tuân thủ chuẩn PEP 8 và mẫu định dạng văn bản khoa học LaTeX.
        \item Hỗ trợ cấu trúc các bài kiểm thử tự động (\texttt{pytest}) cho các module tiền xử lý dữ liệu.
    \end{itemize}
    \item \textbf{Cam đoan về tính xác thực của dữ liệu và kết quả thực nghiệm:}
    \begin{itemize}
        \item 100\% dữ liệu phân tích là bộ dữ liệu thực tế \textit{Hotel Booking Demand} của Jesse Mostipak. Tuyệt đối không sử dụng dữ liệu tự tạo giả định.
        \item Toàn bộ bảng số liệu đánh giá, ma trận nhầm lẫn và các đồ thị được sinh tự động và trung thực từ quá trình thực thi mã nguồn trên máy tính thực tế.
        \item Các diễn giải về nghiệp vụ lưu trú, phân tích tương quan và biện luận nguyên nhân lỗi được các thành viên trong nhóm trực tiếp nghiên cứu và đúc kết.
    \end{itemize}
\end{enumerate}

Nhóm sinh viên hoàn toàn chịu trách nhiệm về nội dung báo cáo và mã nguồn của dự án trước giảng viên hướng dẫn và hội đồng bảo vệ học phần.
